%% ============================================================
%% MSDNet: A Multi-Scale Disentangled Network for Simultaneous
%% Diabetic and Hypertensive Retinopathy Detection
%% Target: Computers in Biology and Medicine (Elsevier, Scopus Q1)
%% Template: elsarticle
%%
%% PLACEHOLDER SYSTEM — search for \RESULT before submission:
%%   \RESULT{label} renders as \textbf{[--]} — fill with real number
%%   \RREF{note}    marks a citation detail to verify
%%   % VERIFY:      comment flagging something to double-check
%% ============================================================

\documentclass[review,12pt]{elsarticle}

%% --- Core packages ---
\usepackage{lineno}          % Required: line numbers for review
\usepackage{amsmath,amssymb,amsfonts}
\usepackage{graphicx}
\usepackage{booktabs}
\usepackage{multirow}
\usepackage{array}
\usepackage{algorithm}
\usepackage{algpseudocode}
\usepackage{xcolor}
\usepackage{hyperref}
\usepackage{microtype}
\usepackage{siunitx}
\usepackage{threeparttable}
\usepackage{caption}
\usepackage{subcaption}

%% --- Placeholder command: fill these after experiments ---
\newcommand{\RESULT}[1]{\textbf{[--]}}
%% Usage: \RESULT{dr_qwk_baseline} → renders as [--] until filled
%% After experiments, do a search-replace: \RESULT{dr_qwk_baseline} → 0.83

%% --- Journal metadata ---
\journal{Computers in Biology and Medicine}

%% --- Line numbers (remove before final submission, keep for review) ---
\linenumbers

\begin{document}

% ---------------------------------------------------------------
% HIGHLIGHTS (Elsevier requirement — 3 to 5 bullet points)
% ---------------------------------------------------------------
\begin{highlights}
\item First unified deep learning framework for simultaneous diabetic and hypertensive retinopathy detection from a single retinal fundus image.
\item Novel contrastive disentanglement loss separates disease-specific feature representations despite shared pathological features (retinal hemorrhages).
\item Feature Pyramid Network enables multi-scale detection of lesions spanning two orders of magnitude in size (1--5 px microaneurysms to 200 px optic disc).
\item Monte Carlo Dropout provides calibrated uncertainty estimates enabling safe automated referral of ambiguous cases to specialist review.
\item Three-stage clinical pipeline integrates patient triage, explainable multi-disease diagnosis, and multilingual patient communication for low-resource deployment.
\end{highlights}

% ---------------------------------------------------------------
% FRONTMATTER
% ---------------------------------------------------------------
\begin{frontmatter}

\title{MSDNet: A Multi-Scale Disentangled Network for Simultaneous
Diabetic and Hypertensive Retinopathy Detection with Uncertainty
Quantification and Explainable AI}

%% --- Authors (fill before submission) ---
\author[lpu]{Shivendra Pratap\corref{cor1}}
\ead{your.email@lpu.in}
\author[lpu]{Supervisor Name}

\cortext[cor1]{Corresponding author.}

\affiliation[lpu]{
  organization={School of Computer Science and Engineering},
  addressline={Lovely Professional University},
  city={Phagwara},
  state={Punjab},
  postcode={144411},
  country={India}
}

% ---------------------------------------------------------------
% ABSTRACT
% ---------------------------------------------------------------
\begin{abstract}
Diabetic retinopathy (DR) and hypertensive retinopathy (HR) are the two
leading causes of preventable blindness in India, affecting over 77 million
diabetic patients of whom fewer than 30\% receive regular retinal screening.
Existing automated diagnostic systems address these conditions in isolation,
despite their frequent co-occurrence and visually overlapping pathological
features — most critically, retinal hemorrhages that are produced by both
diseases independently. This paper presents MSDNet (Multi-Scale Disentangled
Network), a unified multi-task architecture for simultaneous DR grading
and HR detection from a single retinal fundus image, embedded within a
three-stage clinical deployment pipeline.

MSDNet employs an EfficientNet-B0 backbone with a Feature Pyramid Network
(FPN) to capture retinal lesions across two orders of magnitude in size.
Disease-specific classification branches, each equipped with a Convolutional
Block Attention Module (CBAM), enable independent spatial and channel
attention patterns for each condition. An auxiliary vessel segmentation
decoder enforces vascular anatomical understanding in the shared encoder.
The core technical contribution is a novel contrastive disentanglement loss
that penalises cosine similarity between disease-specific feature embeddings,
with clinically motivated margins for pure single-disease and co-occurring
cases. Monte Carlo Dropout with 30 stochastic forward passes provides
calibrated uncertainty estimates, flagging ambiguous cases for specialist
review. Per-branch Grad-CAM++ generates separate explainability heatmaps
per disease class.

The complete system, IRDAS (Integrated Retinal Disease Analysis System),
integrates a XGBoost-based pre-screening triage module using six electronic
health record variables and a multilingual patient report generator supporting
nine Indic languages. Experiments on APTOS 2019, HRDC 2023, IDRiD, and DRIVE
datasets demonstrate a DR quadratic weighted kappa of \RESULT{dr_qwk_full}
on APTOS validation, cross-dataset generalisation kappa of
\RESULT{dr_qwk_idrid} on IDRiD, HR AUC of \RESULT{hr_auc_full}, and
expected calibration error of \RESULT{ece_full}. Ablation studies confirm
the independent contribution of each architectural component.
\end{abstract}

\begin{keyword}
Diabetic Retinopathy \sep
Hypertensive Retinopathy \sep
Multi-Task Learning \sep
Feature Pyramid Network \sep
Contrastive Disentanglement \sep
Explainable AI \sep
Uncertainty Quantification \sep
Retinal Fundus Analysis \sep
Low-Resource Clinical Deployment
\end{keyword}

\end{frontmatter}

% ---------------------------------------------------------------
% 1. INTRODUCTION
% ---------------------------------------------------------------
\section{Introduction}
\label{sec:intro}

Diabetic retinopathy (DR) and hypertensive retinopathy (HR) represent two
of the most prevalent, yet preventable, causes of irreversible vision loss
globally. The International Diabetes Federation (IDF) Diabetes Atlas reports
that approximately 537 million adults were living with diabetes in 2021, a
figure projected to reach 783 million by 2045~\cite{idf2021}. India alone
accounts for 77 million diabetic patients, ranking second globally~\cite{idf2021}.
A systematic review and meta-analysis published in \textit{The Lancet Diabetes
\& Endocrinology} estimated that 22.3\% of people with diabetes have some form
of DR, translating to approximately 103 million individuals
worldwide~\cite{teo2023}. Concurrently, hypertension — the primary driver of
HR — affects an estimated one billion adults globally, with prevalence exceeding
30\% in low- and middle-income countries~\cite{who2023hypertension}.
%% VERIFY: WHO 2023 hypertension figure — check WHO Global Status Report

Both conditions are characterised by a clinically silent early phase in which
significant retinal damage accumulates before any symptomatic visual
disturbance. Among diabetic patients in India, fewer than 30\% receive the
annual retinal screening recommended by clinical guidelines, owing to a severe
shortage of ophthalmic specialists — estimated at 18 ophthalmologists per
million population in rural areas — and the absence of retinal imaging
infrastructure in primary care
settings~\cite{raman2019,npcb2022}.
%% VERIFY: raman2019 — Raman R et al., Diabetic Medicine, India screening statistics
%% VERIFY: npcb2022 — National Programme for Control of Blindness, India

Deep learning has demonstrated clinician-level accuracy for automated DR
detection, most prominently in the landmark study by Gulshan et al., who
reported a sensitivity of 97.5\% and specificity of 93.4\% using an
Inception-V3 architecture on a dataset of 128,175 validated retinal
images~\cite{gulshan2016}. More recently, FDA-approved autonomous screening
systems such as IDx-DR have achieved regulatory clearance for DR screening
without requiring specialist oversight~\cite{abramoff2018}. However, despite
this substantial progress, the existing automated literature presents two
fundamental and unresolved limitations.

\textbf{First}, DR and HR are overwhelmingly modelled as isolated diagnostic
tasks, despite their high clinical co-occurrence. Patients with
poorly-controlled Type 2 diabetes and concurrent hypertension — a highly
prevalent combination in the Indian subcontinent — may exhibit lesions of both
conditions simultaneously in the same retinal image. Running two independent
models doubles computational cost and ignores the joint information that
could improve diagnostic accuracy for both conditions.

\textbf{Second}, both DR and HR produce visually similar pathological signs,
most critically retinal hemorrhages: in DR, hemorrhages arise from capillary
wall breakdown due to hyperglycaemia, while in HR they result from
arteriovenous nicking and increased vascular pressure. A model that learns
shared representations for hemorrhagic features will systematically fail to
attribute the correct aetiology — a diagnostically critical distinction that
determines whether the patient is primarily referred to endocrinology or
nephrology/cardiology.

This paper presents \textbf{MSDNet} (Multi-Scale Disentangled Network), a
unified deep learning architecture that simultaneously grades DR severity
(five classes: grades 0--4 per the International Clinical Diabetic Retinopathy
scale) and detects HR presence (binary classification) from a single retinal
fundus image. MSDNet is embedded within the complete \textbf{IRDAS} (Integrated
Retinal Disease Analysis System) pipeline, which addresses the full clinical
pathway from patient triage to diagnosis to post-diagnosis communication.

The principal contributions of this work are as follows:

\begin{enumerate}
  \item \textbf{Multi-scale unified architecture:} A shared EfficientNet-B0
  encoder with a Feature Pyramid Network extracts retinal features across
  three spatial scales, capturing lesions ranging from single-pixel
  microaneurysms to the full optic disc structure. Disease-specific
  Convolutional Block Attention Module (CBAM) branches provide independent
  spatial and channel attention for each condition, and an auxiliary vessel
  segmentation decoder enforces vascular anatomical understanding.

  \item \textbf{Novel contrastive disentanglement loss:} A clinically motivated
  hinge loss on the cosine similarity between DR and HR branch feature
  embeddings, with separate margins for pure single-disease and co-occurring
  cases. This is the first application of inter-disease contrastive
  disentanglement to simultaneous retinal disease grading.

  \item \textbf{Calibrated uncertainty estimation:} Monte Carlo Dropout with 30
  stochastic forward passes provides per-prediction uncertainty scores. The
  uncertainty threshold for specialist referral is determined via a
  sensitivity-specificity analysis on a held-out validation set.

  \item \textbf{Per-branch explainability:} Grad-CAM++ applied independently
  to each disease branch's spatial attention layer, producing two distinct
  heatmaps per image highlighting the retinal regions contributing to each
  diagnosis.

  \item \textbf{End-to-end clinical pipeline:} A three-stage system
  incorporating XGBoost-based patient triage from electronic health records
  (EHR), the MSDNet diagnostic core, and a multilingual patient report
  generator supporting nine Indic languages, designed for deployment in
  low-resource primary care settings in India.
\end{enumerate}

The remainder of this paper is organised as follows. Section~\ref{sec:related}
reviews related work across automated retinal disease detection, multi-task
learning, attention mechanisms, feature disentanglement, and uncertainty
quantification. Section~\ref{sec:method} describes the complete IRDAS
methodology including the MSDNet architecture and all three pipeline stages.
Section~\ref{sec:datasets} characterises the datasets used.
Section~\ref{sec:exp} details the experimental setup. Section~\ref{sec:results}
presents quantitative and qualitative results. Section~\ref{sec:discussion}
discusses the implications, failure modes, and limitations.
Section~\ref{sec:conclusion} concludes the paper.

% ---------------------------------------------------------------
% 2. RELATED WORK
% ---------------------------------------------------------------
\section{Related Work}
\label{sec:related}

\subsection{Automated Diabetic Retinopathy Detection}

Automated DR detection from retinal fundus images has a substantial research
history predating deep learning. Early systems relied on handcrafted features
and classical machine learning for lesion segmentation and
classification~\cite{niemeijer2007}. The emergence of deep convolutional neural
networks (CNNs) transformed the field: Gulshan et al.~\cite{gulshan2016}
demonstrated that a fine-tuned Inception-V3 network could match the sensitivity
and specificity of board-certified ophthalmologists on a large, multi-ethnic
dataset. Subsequent work established EfficientNet architectures as the
dominant paradigm for DR grading due to their superior parameter
efficiency~\cite{tan2019}. Gargeya and Leng~\cite{gargeya2017} demonstrated
that CNNs could identify diabetic retinopathy using only image-level labels
without requiring pixel-level lesion annotations.

Despite these advances, virtually all published DR systems address grading as
a single-disease task, without modelling the clinical context of co-morbid
conditions. Furthermore, most studies report performance on clean benchmark
datasets that do not reflect the heterogeneous image quality encountered in
primary care clinics in developing countries~\cite{raman2019}.

\subsection{Automated Hypertensive Retinopathy Detection}

Hypertensive retinopathy manifests as a spectrum of vascular changes including
arteriovenous nicking, silver and copper wiring of arterioles, flame
haemorrhages, and papilloedema in severe cases. The Keith-Wagener-Barker
classification~\cite{kwb1939} provides the standard clinical grading scheme.
Compared to DR, the automated detection literature for HR is substantially
smaller. Wong and Mitchell~\cite{wong2004} characterised the epidemiology of
hypertensive retinopathy in population-based studies. More recent deep learning
approaches~\cite{dai2021} have applied CNN classifiers to fundus images for
HR grading, but published systems remain limited in number, and no prior work
has attempted joint detection of HR alongside DR in a unified model.
The 2023 Hypertensive Retinopathy Detection Challenge (HRDC) provided the
first publicly available large-scale labelled dataset for HR automated
grading~\cite{hrdc2023}, which this work utilises.
%% VERIFY: dai2021 — Dai L et al., paper on HR detection with CNNs
%% VERIFY: hrdc2023 — HRDC 2023 challenge paper reference

\subsection{Multi-Task Learning in Medical Imaging}

Multi-task learning (MTL) improves generalisation by sharing representations
across related tasks~\cite{crawshaw2020}. In retinal analysis, MTL has been
applied to simultaneous segmentation and classification of DR
lesions~\cite{li2019}, as well as joint optic disc and cup
segmentation for glaucoma~\cite{fu2018}.
%% VERIFY: fu2018 — Fu H et al., joint optic disc cup segmentation, MICCAI

A fundamental challenge in multi-task retinal analysis is shared visual
features across pathologically distinct diseases. Li et al.~\cite{li2019}
addressed intra-disease MTL (e.g., lesion segmentation and severity grading
for a single disease) but did not address inter-disease feature entanglement.
No prior work has proposed a loss function that explicitly disentangles
representations of two co-occurring retinal diseases within a unified
architecture. This gap motivates the contrastive disentanglement loss proposed
in this paper.

\subsection{Attention Mechanisms for Retinal Feature Learning}

Attention mechanisms enable CNNs to emphasise pathologically relevant regions
while suppressing background noise. The Squeeze-and-Excitation (SE)
block~\cite{hu2018} provides channel-wise attention; the Convolutional Block
Attention Module (CBAM)~\cite{woo2018} extends this with complementary spatial
attention. In retinal imaging, attention mechanisms have been applied to
focus on haemorrhages and exudates in DR detection and to optic cup boundaries
in glaucoma screening. Our work advances this by assigning \textit{independent}
CBAM modules to each disease branch, enabling each to learn entirely distinct
spatial and channel attention patterns, rather than sharing a single attention
module across tasks.
%% VERIFY: hu2018 — Hu J et al., Squeeze-and-Excitation Networks, CVPR 2018

\subsection{Feature Disentanglement}

Feature disentanglement aims to learn representations where different factors
of variation are encoded in separable subspaces. In generative modelling,
$\beta$-VAE~\cite{higgins2017} enforces disentangled latent spaces via an
information-theoretic penalty.
%% VERIFY: higgins2017 — Higgins I et al., beta-VAE, ICLR 2017
In supervised settings, contrastive losses~\cite{khosla2020} have been used
to enforce class-separable representations. The application of contrastive
disentanglement to simultaneously detected, co-occurring medical conditions
is, to the best of our knowledge, novel. The clinical motivation differs
fundamentally from standard contrastive learning: rather than separating
healthy from pathological samples, we separate disease-specific representations
within pathological samples, guided by the known medical difference in
aetiology between DR and HR haemorrhages.

\subsection{Multi-Scale Feature Extraction}

Retinal pathologies span a wide range of spatial scales: microaneurysms are
approximately 10--100 $\mu$m in diameter (corresponding to 1--5 pixels in
standard fundus images), while the optic disc subtends approximately 1,500 $\mu$m
(100--200 pixels). Standard CNNs with fixed receptive fields perform poorly on
this multi-scale problem. The Feature Pyramid Network (FPN)~\cite{lin2017fpn}
addresses this by fusing feature maps across multiple backbone stages in a
top-down pathway with lateral connections, providing rich, multi-scale
representations without a substantial computational overhead. FPN has been
applied in retinal vessel segmentation~\cite{meyer2021} but not specifically
to the problem of multi-scale lesion detection across co-occurring
diseases.
%% VERIFY: meyer2021 — a paper applying FPN to retinal vessel segmentation

\subsection{Uncertainty Quantification in Clinical AI}

A central barrier to clinical adoption of AI diagnostic systems is the absence
of calibrated confidence estimates. Standard softmax outputs are known to be
overconfident and poorly calibrated~\cite{guo2017}. Begoli et
al.~\cite{begoli2019} argued that uncertainty quantification is a prerequisite
for the responsible deployment of AI in clinical settings. Gal and
Ghahramani~\cite{gal2016} demonstrated that Monte Carlo Dropout provides a
practical Bayesian approximation, with test-time dropout masks simulating
samples from the posterior distribution over model weights. In retinal imaging,
calibrated uncertainty has been applied to diabetic retinopathy screening to
identify images requiring human grader review~\cite{leibig2017}, but not to
the simultaneous multi-disease scenario with calibration evaluated per disease
branch.
%% VERIFY: guo2017 — Guo C et al., On Calibration of Modern Neural Networks, ICML 2017
%% VERIFY: leibig2017 — Leibig C et al., leveraging uncertainty in DR screening

\subsection{Explainability in Retinal AI}

Clinical acceptance of AI diagnostic tools requires interpretable evidence for
each prediction. Grad-CAM~\cite{selvaraju2017} provides class activation maps
by weighting feature map channels by the gradient of the target class score.
Grad-CAM++~\cite{chattopadhyay2018} extends this using second-order gradients,
producing sharper localisation of multiple instances of the same class —
critical for retinal images with scattered microaneurysms. Prior XAI work in
retinal AI generates a single heatmap per image~\cite{sayres2019}. Our
approach generates \textit{separate} heatmaps for the DR and HR branches,
enabling clinicians to visually confirm that each branch attends to its
disease-specific lesions — a qualitative validation of the disentanglement
mechanism that single-branch systems cannot provide.
%% VERIFY: sayres2019 — Sayres R et al., using XAI in DR screening

% ---------------------------------------------------------------
% 3. METHODOLOGY
% ---------------------------------------------------------------
\section{Methodology}
\label{sec:method}

\subsection{System Overview}

IRDAS (Integrated Retinal Disease Analysis System) operates as a three-stage
clinical pipeline, illustrated in Fig.~\ref{fig:overview}. \textbf{Stage~1}
is a pre-screening triage module that takes six clinical biomarkers from a
patient's electronic health record and outputs a ranked risk score for
undetected DR, prioritising which patients should receive fundus imaging.
\textbf{Stage~2} is the core MSDNet model, which processes a retinal fundus
image and outputs DR grade, HR prediction, per-disease uncertainty scores,
and two Grad-CAM++ heatmaps. \textbf{Stage~3} is a multilingual communication
module that synthesises the Stage~2 outputs into a plain-language patient
report in the detected Indic language of the patient.

\begin{figure}[htbp]
  \centering
  %% FIGURE PLACEHOLDER — replace with actual architecture diagram
  %% The figure should show:
  %%   (a) Three-stage pipeline overview (horizontal flow diagram)
  %%   (b) Detailed MSDNet architecture: EfficientNet blocks → FPN →
  %%       Split into DR branch + HR branch + Vessel decoder
  %%       with CBAM detail inset, loss terms annotated
  \framebox[0.95\columnwidth]{\parbox{0.9\columnwidth}{\centering\vspace{2cm}
  [Figure 1: IRDAS system overview and MSDNet architecture diagram.
   Generate from draw.io or matplotlib after implementation.]\vspace{2cm}}}
  \caption{Overview of IRDAS and the MSDNet architecture. (a) The three-stage
  pipeline: pre-screening triage from EHR data (Stage 1), multi-disease retinal
  image analysis with MSDNet (Stage 2), and multilingual patient report
  generation (Stage 3). (b) Detailed MSDNet architecture showing the shared
  EfficientNet-B0 encoder, Feature Pyramid Network, disease-specific CBAM
  branches with MC Dropout, auxiliary vessel segmentation decoder, and
  per-branch Grad-CAM++ explainability.}
  \label{fig:overview}
\end{figure}

\subsection{Preprocessing Pipeline}

A standardised five-step preprocessing pipeline is applied uniformly to all
fundus images regardless of source dataset. A visual comparison of the
preprocessing stages is shown in Fig.~\ref{fig:preprocessing}.

\begin{enumerate}
  \item \textbf{Ben Graham normalisation~\cite{graham2015}:} Removes uneven
  illumination by subtracting a Gaussian-blurred version of the image:
  \begin{equation}
    I_{\text{out}} = \alpha I - \alpha G_{\sigma}(I) + \beta,
    \label{eq:ben_graham}
  \end{equation}
  where $\alpha = 4$, $\sigma = \text{image\_width}/30$, and $\beta = 128$
  centres the intensity distribution. This step normalises the uneven
  vignetting characteristic of fundus cameras.

  \item \textbf{CLAHE~\cite{zuiderveld1994}:} Contrast Limited Adaptive
  Histogram Equalization applied to the L channel in LAB colour space with clip
  limit 2.0 and $8 \times 8$ tile grid. Green channel analysis is preferred as
  retinal microvasculature achieves maximum contrast against the red fundus
  background in this channel.

  \item \textbf{Optic disc suppression:} The optic disc is localised as the
  maximum-intensity circular region using the Hough Circle Transform and
  suppressed via Navier-Stokes inpainting. This prevents the model from
  misclassifying the bright optic disc as a cluster of hard exudates — a
  documented failure mode of models without explicit OD handling.

  \item \textbf{Circular crop:} The black border surrounding the circular
  fundus field of view is removed using contour-based bounding box extraction,
  removing uninformative pixels that could otherwise consume receptive field
  capacity.

  \item \textbf{Resize and normalise:} Images are resized to $224 \times 224$
  pixels using bicubic interpolation and normalised using ImageNet statistics
  ($\mu = [0.485, 0.456, 0.406]$, $\sigma = [0.229, 0.224, 0.225]$).
\end{enumerate}

\begin{figure}[htbp]
  \centering
  %% FIGURE PLACEHOLDER — side-by-side comparison:
  %% Raw image | After Ben Graham | After CLAHE | After OD suppression | Final
  %% Use a Grade 3 APTOS image for maximum visual impact
  \framebox[0.95\columnwidth]{\parbox{0.9\columnwidth}{\centering\vspace{1.5cm}
  [Figure 2: Preprocessing pipeline visualisation. Five-column image strip
   showing a Grade 3 DR fundus image at each preprocessing stage.]\vspace{1.5cm}}}
  \caption{Effect of the preprocessing pipeline on a representative Grade 3
  diabetic retinopathy fundus image. From left to right: (a) raw image,
  (b) after Ben Graham normalisation, (c) after CLAHE, (d) after optic disc
  suppression, (e) final preprocessed image. Note the enhanced visibility of
  haemorrhage dots and vessel boundaries in (e) compared to (a).}
  \label{fig:preprocessing}
\end{figure}

Training data augmentation is applied only during training and consists of
geometric and photometric transformations, detailed in
Table~\ref{tab:augmentation}.

\begin{table}[htbp]
\caption{Data Augmentation Pipeline (Training Only)}
\label{tab:augmentation}
\centering
\begin{threeparttable}
\begin{tabular}{llc}
\toprule
\textbf{Transformation} & \textbf{Parameters} & \textbf{Probability} \\
\midrule
Horizontal flip       & ---                                   & 0.50 \\
Vertical flip         & ---                                   & 0.50 \\
Random rotation       & $[-30°, +30°]$                        & 0.60 \\
Shift-scale-rotate    & shift 0.1, scale 0.15, rotate 30°     & 0.60 \\
Colour jitter\tnote{a} & brightness 0.2, contrast 0.2, saturation 0.1 & 0.50 \\
Gaussian noise        & $\sigma \in [0, 0.05]$                & 0.30 \\
Gaussian blur         & kernel $\in \{3, 5\}$                 & 0.20 \\
Elastic transform     & $\alpha = 120$, $\sigma = 6$          & 0.30 \\
Random brightness/contrast & $\pm 0.3$                        & 0.40 \\
\bottomrule
\end{tabular}
\begin{tablenotes}
\small
\item[a] Hue shift is excluded from colour jitter to preserve diagnostic
colour information (the red-orange hue of haemorrhages vs. the yellow-white
of exudates is a key discriminating feature).
\end{tablenotes}
\end{threeparttable}
\end{table}

\subsection{Multi-Scale Feature Extraction}

The backbone is EfficientNet-B0~\cite{tan2019} with ImageNet-pretrained weights.
EfficientNet-B0 was selected for its compound scaling strategy (simultaneous
scaling of depth, width, and resolution) and parameter efficiency (5.3 million
parameters), making it feasible for deployment on standard clinical hardware
without specialised GPU requirements.

Using the \texttt{features\_only} extraction mode, intermediate feature maps
are obtained at three backbone stages, yielding a natural feature pyramid for
a $224 \times 224$ input:

\begin{itemize}
  \item $\mathbf{P}_3 \in \mathbb{R}^{B \times 40 \times 28 \times 28}$:
  stride 8, fine-grained features capturing microaneurysms (1--5 px), small
  haemorrhage dots, and vessel edge details.
  \item $\mathbf{P}_4 \in \mathbb{R}^{B \times 112 \times 14 \times 14}$:
  stride 16, medium-scale features capturing haemorrhage clusters, exudate
  patches, and localised vessel abnormalities.
  \item $\mathbf{P}_5 \in \mathbb{R}^{B \times 320 \times 7 \times 7}$:
  stride 32, coarse features capturing the optic disc, global vessel
  calibre patterns, and arteriovenous relationships.
\end{itemize}

The clinical significance of multi-scale extraction is that no single
receptive field size can simultaneously resolve a 2-pixel microaneurysm and
the 200-pixel optic disc. Standard single-scale classifiers must implicitly
choose between fine and coarse features, creating an architectural
bias~\cite{lin2017fpn}.

\subsection{Feature Pyramid Network}

Multi-scale features are fused using a Feature Pyramid Network
(FPN)~\cite{lin2017fpn}. Lateral $1 \times 1$ convolutions unify channel
dimensions across pyramid levels, and a top-down pathway progressively
fuses coarse semantic information with fine spatial detail:

\begin{equation}
  \mathbf{P}'_5 = \text{Conv}_{1\times1}(\mathbf{P}_5),
  \label{eq:fpn1}
\end{equation}
\begin{equation}
  \mathbf{P}'_4 = \text{Conv}_{1\times1}(\mathbf{P}_4) +
  \text{Up}(\mathbf{P}'_5),
  \label{eq:fpn2}
\end{equation}
\begin{equation}
  \mathbf{P}'_3 = \text{Conv}_{1\times1}(\mathbf{P}_3) +
  \text{Up}(\mathbf{P}'_4),
  \label{eq:fpn3}
\end{equation}
\begin{equation}
  \mathbf{F}_{\text{fpn}} = \text{BN}\!\left(\text{Conv}_{3\times3}(\mathbf{P}'_3)\right)
  \in \mathbb{R}^{B \times 256 \times 28 \times 28},
  \label{eq:fpn4}
\end{equation}

where $\text{Up}(\cdot)$ denotes nearest-neighbour upsampling, and batch
normalisation (BN) followed by a $3 \times 3$ convolution smooths artefacts
from upsampling. The output $\mathbf{F}_{\text{fpn}}$ is the unified
multi-scale feature map that serves as input to both disease branches.

\subsection{Disease-Specific Classification Branches}

To prevent feature entanglement between DR and HR representations, each
disease is assigned a dedicated classification branch with its own CBAM
attention module. Let $d \in \{\text{DR}, \text{HR}\}$ index the branches.
For each branch:

\begin{equation}
  \tilde{\mathbf{x}}_d = \text{CBAM}_d(\mathbf{F}_{\text{fpn}}),
  \label{eq:cbam}
\end{equation}
\begin{equation}
  \mathbf{f}_d = \text{Dropout}_p\!\left(\text{GAP}(\tilde{\mathbf{x}}_d)\right)
  \in \mathbb{R}^{B \times 256},
  \label{eq:branch_feat}
\end{equation}
\begin{equation}
  \hat{\mathbf{y}}_d = W_d \mathbf{f}_d + \mathbf{b}_d,
  \label{eq:branch_head}
\end{equation}

where $\text{CBAM}_d$ applies sequential channel attention (reduction ratio 16,
shared MLP) and spatial attention ($7 \times 7$ convolution); GAP is global
average pooling; Dropout is applied with $p = 0.3$; and $W_d$, $\mathbf{b}_d$
are the branch-specific linear layer weights. The DR branch produces a
five-dimensional logit vector (grades 0--4), and the HR branch produces a
scalar logit for binary classification. Both branches return the 256-dimensional
embedding $\mathbf{f}_d$, which is used in the contrastive disentanglement
loss.

CBAM applies two sequential submodules. The channel attention submodule
\cite{woo2018} recalibrates channel-wise feature responses:

\begin{equation}
  \mathbf{M}_c(\mathbf{F}) = \sigma\!\left(\text{MLP}\!\left(\text{AvgPool}(\mathbf{F})\right)
  + \text{MLP}\!\left(\text{MaxPool}(\mathbf{F})\right)\right),
  \label{eq:cbam_channel}
\end{equation}

and the spatial attention submodule identifies where to focus within the
feature map:

\begin{equation}
  \mathbf{M}_s(\mathbf{F}) = \sigma\!\left(\text{Conv}_{7\times7}\!\left(
  [\text{AvgPool}_c(\mathbf{F});\text{MaxPool}_c(\mathbf{F})]\right)\right),
  \label{eq:cbam_spatial}
\end{equation}

where $[\cdot;\cdot]$ denotes channel concatenation and $\text{Pool}_c$
denotes pooling along the channel dimension. The spatial attention convolution
layer of each branch's CBAM ($\text{Conv}_{7\times7}$ in
Eq.~\ref{eq:cbam_spatial}) is used as the Grad-CAM++ target layer.

\subsection{Auxiliary Vessel Segmentation Decoder}

An auxiliary U-Net-style decoder~\cite{ronneberger2015} is attached directly
to the backbone feature pyramid ($\mathbf{P}_3$, $\mathbf{P}_4$, $\mathbf{P}_5$),
parallel to but independent of the FPN-disease branch pathway. The decoder
consists of three upsampling blocks:

\begin{align}
  \mathbf{D}_5 &= \text{UpBlock}(\mathbf{P}_5), \quad \mathbf{D}_5 \in
  \mathbb{R}^{B \times 112 \times 14 \times 14} \label{eq:dec1}\\
  \mathbf{D}_4 &= \text{UpBlock}(\text{Cat}[\mathbf{D}_5;\mathbf{P}_4]), \quad
  \mathbf{D}_4 \in \mathbb{R}^{B \times 64 \times 28 \times 28} \label{eq:dec2}\\
  \mathbf{D}_3 &= \text{UpBlock}(\text{Cat}[\mathbf{D}_4;\mathbf{P}_3]), \quad
  \mathbf{D}_3 \in \mathbb{R}^{B \times 32 \times 56 \times 56} \label{eq:dec3}\\
  V_{\text{pred}} &= \sigma\!\left(\text{Conv}_{1\times1}(\text{Up}(\mathbf{D}_3))\right)
  \in \mathbb{R}^{B \times 1 \times 224 \times 224}, \label{eq:dec4}
\end{align}

where each UpBlock consists of two $3 \times 3$ convolution-BN-ReLU operations
followed by bilinear upsampling, and $\text{Cat}[\cdot;\cdot]$ denotes channel
concatenation implementing the U-Net skip connection. The decoder adds
approximately 0.8 million parameters.

The vessel segmentation auxiliary task is domain-justified: both DR and HR
are fundamentally diseases of the retinal vasculature. By requiring the shared
encoder to produce features from which vessel topology can be recovered, we
enforce that the encoder learns representations that are \textit{anatomically
grounded} rather than dependent on superficial texture cues. This regularisation
effect on the backbone is the primary purpose of the auxiliary task; vessel
segmentation performance, while reported, is a secondary goal.

The decoder is trained on the DRIVE dataset (Section~\ref{sec:datasets}) in
addition to CHASE\_DB1 and STARE to augment the limited number of vessel
segmentation training images. The decoder is active only during training and
is entirely removed at inference, adding zero latency overhead to deployment.

\subsection{Training Loss Functions}

The total training objective combines four loss terms:

\begin{equation}
  \mathcal{L}_{\text{total}} = \mathcal{L}_{\text{DR}} + \mathcal{L}_{\text{HR}}
  + \lambda_v \mathcal{L}_{\text{vessel}} + \lambda_c \mathcal{L}_{\text{dis}},
  \label{eq:total_loss}
\end{equation}

with loss weights $\lambda_v = 0.5$ and $\lambda_c = 0.3$. Sensitivity
analysis of these weights is reported in Section~\ref{sec:results}.

\subsubsection{Focal Loss for DR Grading}

The APTOS 2019 dataset exhibits severe class imbalance, with Grade 0
(no DR) comprising approximately 49\% of training samples and Grade 4
(proliferative DR) comprising 9\%. Standard cross-entropy loss would bias the
model towards the majority class. We apply focal loss~\cite{lin2017focal} to
downweight the contribution of well-classified easy examples:

\begin{equation}
  \mathcal{L}_{\text{DR}} = -\sum_{i=1}^{N} (1 - p_{t,i})^{\gamma} \log(p_{t,i}),
  \quad \gamma = 2.0,
  \label{eq:focal}
\end{equation}

where $p_t$ is the model probability for the ground-truth class and $\gamma$
is the focusing parameter that reduces the loss for high-confidence predictions.
Class weights inversely proportional to class frequency are additionally applied.

\subsubsection{Binary Cross-Entropy for HR Detection}

\begin{equation}
  \mathcal{L}_{\text{HR}} = -\left[y \log \hat{y} + (1-y)\log(1-\hat{y})\right],
  \label{eq:bce}
\end{equation}

where $\hat{y} = \sigma(\hat{y}_{\text{HR}})$ is the sigmoid-activated HR
branch output.

\subsubsection{Dice-BCE Loss for Vessel Segmentation}

\begin{equation}
  \mathcal{L}_{\text{vessel}} = 1 - \frac{2|V_{\text{pred}} \cap V_{\text{gt}}|
  + \varepsilon}{|V_{\text{pred}}| + |V_{\text{gt}}| + \varepsilon}
  + \text{BCE}(V_{\text{pred}}, V_{\text{gt}}),
  \label{eq:vessel}
\end{equation}

where $\varepsilon = 1$ is a smoothing constant preventing division by zero.

\subsubsection{Contrastive Disentanglement Loss (Core Contribution)}

The fundamental clinical observation motivating this loss is that DR and HR
both cause retinal haemorrhages, but through distinct pathophysiological
mechanisms: DR haemorrhages arise from capillary wall breakdown due to
hyperglycaemia-induced pericyte apoptosis, while HR haemorrhages result from
arteriovenous nicking and elevated hydrostatic pressure. A model without
explicit disentanglement will learn a single shared representation for
``haemorrhage'', preventing it from identifying which disease produced any
given haemorrhage in a co-occurring case.

We define the inter-disease cosine similarity between branch embeddings:

\begin{equation}
  s_i = \frac{\mathbf{f}_{\text{DR},i} \cdot \mathbf{f}_{\text{HR},i}}
  {\|\mathbf{f}_{\text{DR},i}\|_2 \, \|\mathbf{f}_{\text{HR},i}\|_2},
  \label{eq:cosine}
\end{equation}

where $\mathbf{f}_{\text{DR},i}$ and $\mathbf{f}_{\text{HR},i}$ are the
$\ell_2$-normalised branch embeddings for sample $i$. Three clinically
distinct cases receive different treatment:

\begin{description}
  \item[Pure single-disease samples ($\mathcal{P}$):] Images with DR only
  (HR absent) or HR only (DR absent). The two branches should attend to
  entirely distinct retinal features. A strict margin $m_p = 0.1$ is applied.

  \item[Co-occurring samples ($\mathcal{C}$):] Images with both DR and HR
  present. Some feature sharing is clinically appropriate (the same vessels
  are affected by both conditions), but representations must remain
  sufficiently separated to support independent diagnosis. A relaxed margin
  $m_c = 0.3$ is applied.

  \item[Healthy samples:] Neither DR nor HR present. No disentanglement
  constraint is applied, as neither branch has a positive signal to
  disentangle.
\end{description}

The total disentanglement loss is:

\begin{equation}
  \mathcal{L}_{\text{dis}} =
  \frac{\sum_{i \in \mathcal{P}} [s_i - m_p]_+}{|\mathcal{P}|}
  + \frac{\sum_{j \in \mathcal{C}} [s_j - m_c]_+}{|\mathcal{C}|},
  \label{eq:dis_loss}
\end{equation}

where $[\cdot]_+ = \max(0, \cdot)$ is the hinge function, and $|\mathcal{P}|$,
$|\mathcal{C}|$ are the number of samples in each set within the batch,
clamped to a minimum of 1 to avoid division by zero. Embeddings are
$\ell_2$-normalised before computing $s_i$ to ensure cosine similarity
is bounded in $[-1, 1]$ regardless of embedding magnitude.

\subsection{Multi-Dataset Joint Training Strategy}

A critical implementation challenge arises from the disjoint label spaces of
the training datasets: APTOS 2019 provides only DR grade labels (no HR labels),
and HRDC 2023 provides only HR labels (no DR grade labels). Standard multi-task
training on a single combined batch is therefore not directly applicable.

We adopt an \textit{alternating batch} training strategy, detailed in
Algorithm~\ref{alg:training}. On each training iteration, a batch is drawn
from either APTOS (with probability $q$) or HRDC (with probability $1-q$).
Loss terms are masked based on which dataset the batch originates from:
APTOS batches compute only $\mathcal{L}_{\text{DR}}$ and
$\mathcal{L}_{\text{vessel}}$; HRDC batches compute only $\mathcal{L}_{\text{HR}}$.
The contrastive disentanglement loss $\mathcal{L}_{\text{dis}}$ is computed
only on batches that contain co-occurring or pure single-disease samples with
both labels defined, using the HRDC validation annotations for HR labels on
a shared image subset.

\begin{algorithm}[htbp]
\caption{MSDNet Multi-Dataset Alternating Training}
\label{alg:training}
\begin{algorithmic}[1]
\Require APTOS DataLoader $\mathcal{D}_A$, HRDC DataLoader $\mathcal{D}_H$,
DRIVE DataLoader $\mathcal{D}_V$, sampling probability $q=0.6$
\Require Trained model $\theta$, optimiser, loss weights
$\lambda_v$, $\lambda_c$
\State Initialise $\theta$ from ImageNet pretrained EfficientNet-B0 weights
\For{each epoch}
  \State Create iterators $\text{iter}_A$, $\text{iter}_H$, $\text{iter}_V$
  \For{each training step}
    \State Draw $u \sim \text{Uniform}(0,1)$
    \If{$u < q$} \Comment{Sample from APTOS}
      \State $(x, y_{\text{DR}}) \leftarrow \text{next}(\text{iter}_A)$
      \State $\hat{y} \leftarrow \text{MSDNet}(x)$
      \State $\mathcal{L} \leftarrow \mathcal{L}_{\text{DR}}(\hat{y}_{\text{DR}},
        y_{\text{DR}})$
    \Else \Comment{Sample from HRDC}
      \State $(x, y_{\text{HR}}) \leftarrow \text{next}(\text{iter}_H)$
      \State $\hat{y} \leftarrow \text{MSDNet}(x)$
      \State $\mathcal{L} \leftarrow \mathcal{L}_{\text{HR}}(\hat{y}_{\text{HR}},
        y_{\text{HR}})$
    \EndIf
    \If{vessel step (every 5th step)}
      \State $(x_v, m_v) \leftarrow \text{next}(\text{iter}_V)$
      \State $\mathcal{L} \leftarrow \mathcal{L} +
        \lambda_v \mathcal{L}_{\text{vessel}}(\text{MSDNet}(x_v), m_v)$
    \EndIf
    \If{both DR and HR labels available in batch}
      \State $\mathcal{L} \leftarrow \mathcal{L} +
        \lambda_c \mathcal{L}_{\text{dis}}(\mathbf{f}_{\text{DR}},
        \mathbf{f}_{\text{HR}}, y_{\text{DR}}, y_{\text{HR}})$
    \EndIf
    \State $\theta \leftarrow \theta - \eta \nabla_\theta \mathcal{L}$
    \Comment{AdamW update}
  \EndFor
  \State Evaluate on validation sets; save checkpoint if improved
\EndFor
\end{algorithmic}
\end{algorithm}

The sampling probability $q = 0.6$ (60\% APTOS, 40\% HRDC) was selected
to reflect the relative dataset sizes. Sensitivity to this hyperparameter
is reported in Section~\ref{sec:ablation}.

\subsection{Uncertainty Estimation via Monte Carlo Dropout}

At inference, dropout layers within both branch heads are kept active (training
mode) and the input image is passed through the network $N$ times with
different dropout masks. For the DR branch:

\begin{equation}
  \bar{\mathbf{p}}_{\text{DR}} = \frac{1}{N}\sum_{t=1}^{N}
  \text{softmax}(\hat{\mathbf{y}}_{\text{DR}}^{(t)}),
  \label{eq:mc_mean}
\end{equation}
\begin{equation}
  u_{\text{DR}} = \frac{1}{K}\sum_{k=1}^{K} \text{Var}_t
  \!\left[\text{softmax}(\hat{\mathbf{y}}_{\text{DR}}^{(t)})_k\right],
  \label{eq:mc_var}
\end{equation}

where $K=5$ is the number of DR classes and the variance is taken over $N=30$
stochastic passes. The choice $N=30$ was determined empirically by observing
the convergence of $u_{\text{DR}}$ as a function of $N$; uncertainty estimates
stabilised at $N \approx 25$ on the validation set, and $N=30$ was selected
with a small margin. The uncertainty score is separately computed for the HR
branch using the predicted sigmoid probability variance.

A prediction is flagged for specialist review when $u_{\text{DR}} >
\tau_{\text{DR}}$ or $u_{\text{HR}} > \tau_{\text{HR}}$. The referral
thresholds $\tau$ are selected post hoc on the validation set by maximising
the F1-score of the referral decision, subject to a minimum sensitivity
constraint of 95\% for Grade $\geq 2$ DR cases (clinically referable).

\subsection{Per-Branch Explainability (Grad-CAM++)}

Grad-CAM++~\cite{chattopadhyay2018} is applied independently to each disease
branch. The target layer for each branch is the spatial attention convolution
in its CBAM module (the $\text{Conv}_{7\times7}$ in Eq.~\ref{eq:cbam_spatial}).
This choice ensures the heatmap reflects not just what the backbone extracted
but what the disease-specific branch chose to attend to — the most clinically
informative attribution. Two heatmaps are generated per inference:
$H_{\text{DR}} \in [0,1]^{224 \times 224}$ highlighting regions contributing
to the DR prediction, and $H_{\text{HR}}$ for HR. The peak activation
quadrant in each heatmap is identified using ophthalmological convention
(superior/inferior $\times$ nasal/temporal $\times$ central macula) and passed
as a text description to the Stage~3 report generation module.

\subsection{Stage 1: Clinical Triage Module}

The triage module addresses the capacity bottleneck upstream of fundus imaging:
in low-resource settings, a primary care clinic may have access to a fundus
camera for only a few hours per week, yet may have hundreds of diabetic
patients registered. The triage module ranks patients by estimated risk of
having undetected referable DR using only six EHR variables, without requiring
a fundus image.

An XGBoost gradient-boosted classifier~\cite{chen2016} is trained on the
NHANES tabular dataset. Input features are: glycated haemoglobin (HbA1c),
systolic blood pressure, age, diabetes duration in years, body mass index
(BMI), and binary insulin usage indicator. The output is a scalar probability
$\hat{r} \in [0,1]$ representing the estimated probability of referable DR
($\geq$ Grade 2). Class imbalance is addressed via the
\texttt{scale\_pos\_weight} parameter, set to the ratio of negative to positive
samples in the training set. SHAP~\cite{lundberg2017} values are computed for
each prediction to provide individual-level feature attribution, enabling the
clinic to understand which risk factors drove a patient's triage score.
%% VERIFY: lundberg2017 — Lundberg SM, Lee SI. SHAP values paper, NIPS 2017

\subsection{Stage 3: Multilingual Patient Communication}

The multilingual communication module converts the structured output of MSDNet
(disease grades, uncertainty flags, and heatmap region descriptions) into a
plain-language patient explanation in the patient's detected Indic language.
The patient's language is identified automatically using the multimodal
language identification model developed by the authors in prior
work~\cite{pratap2026ieee}, which achieves 93.47\% accuracy across 15 Indic
languages.
%% Fill with correct citation for your accepted IEEE paper once published

A language model (Gemini Pro via LangChain) generates a structured explanation
covering: (1) what was found, (2) what it means for the patient's daily life,
and (3) specific next steps (e.g., ``You must visit an eye specialist within
7 days''). The prompt is templated with disease severity, affected retinal
quadrant from the heatmap analysis, and confidence level. Temperature is set
to 0.3 to ensure consistent, conservative medical communication. Report quality
is evaluated using BLEU-4, BERTScore, and a clinician usability rating (1--5)
on a sample of 50 reports evaluated by three ophthalmologists.

% ---------------------------------------------------------------
% 4. DATASETS
% ---------------------------------------------------------------
\section{Datasets}
\label{sec:datasets}

This study uses exclusively publicly available, de-identified benchmark
datasets. No additional patient data was collected. No institutional ethical
review was required. Table~\ref{tab:datasets} summarises all datasets and
their roles.

\begin{table}[htbp]
\caption{Summary of Datasets Used}
\label{tab:datasets}
\centering
\begin{tabular}{lrllll}
\toprule
\textbf{Dataset} & \textbf{Images} & \textbf{Labels} & \textbf{Train}
& \textbf{Val} & \textbf{Test / Role} \\
\midrule
APTOS 2019     & 3,662  & DR grades 0--4 & 2,930 & 366  & 366 (DR primary) \\
HRDC 2023      & 1,200  & HR binary      & 840   & 120  & 240 (HR primary) \\
IDRiD          & 516    & DR grades 0--4 & ---   & ---  & 516 (cross-dataset) \\
DRIVE          & 40     & Vessel masks   & 20    & ---  & 20  (vessel aux.) \\
CHASE\_DB1     & 28     & Vessel masks   & 22    & ---  & 6   (vessel aux.) \\
STARE          & 20     & Vessel masks   & 16    & ---  & 4   (vessel aux.) \\
NHANES         & tabular & DR binary     & 80\%  & 20\% & 5-fold CV (triage) \\
\bottomrule
\end{tabular}
\end{table}

\textbf{APTOS 2019 Blindness Detection:} Released for the 2019 Kaggle
competition, this dataset contains 3,662 retinal fundus images from rural
screening camps in India, graded for DR severity on a five-class scale by
trained graders. The dataset is particularly appropriate for this study as
its images originate from Indian clinical settings. Class distribution is
highly imbalanced: Grade 0 (49.3\%), Grade 1 (10.1\%), Grade 2 (27.3\%),
Grade 3 (5.3\%), Grade 4 (8.1\%). The primary metric is Quadratic Weighted
Kappa (QWK).

\textbf{HRDC 2023:} The Hypertensive Retinopathy Detection Challenge 2023
dataset provides fundus images with HR severity labels following the
Keith-Wagener-Barker~\cite{kwb1939} grading scheme. We convert this to binary
classification (HR present / absent) to focus on detection, which is the
primary clinical need in screening contexts. To our knowledge, this is among
the first papers to use this dataset for a published comparative study.

\textbf{IDRiD (Indian Diabetic Retinopathy Image Dataset):} Collected from
a single Indian eye care centre, IDRiD~\cite{porwal2020} provides 516 fundus
images with expert DR grading and pixel-level lesion segmentation masks. We
use IDRiD \textit{exclusively for cross-dataset testing} — no IDRiD images
are used in training or validation. This evaluates whether MSDNet generalises
to real Indian clinical images, which is the primary deployment target.

\textbf{DRIVE, CHASE\_DB1, and STARE:} Three standard retinal vessel
segmentation benchmark datasets~\cite{staal2004,fraz2012} providing 68 combined
training images with expert vessel segmentation masks, used to train the
auxiliary vessel decoder.
%% VERIFY: fraz2012 — Fraz MM et al., CHASE_DB1 and comparative vessel paper

\textbf{NHANES:} National Health and Nutrition Examination Survey data
(cycles 2005--2008 and 2011--2016) providing tabular health measurements
including HbA1c, blood pressure, BMI, diabetes duration, and insulin usage
alongside ophthalmology examination findings used to derive the DR binary
outcome (referable DR: Grade $\geq 2$).

% ---------------------------------------------------------------
% 5. EXPERIMENTAL SETUP
% ---------------------------------------------------------------
\section{Experimental Setup}
\label{sec:exp}

\textbf{Hardware:} All experiments are conducted on a single NVIDIA Tesla T4
GPU (16 GB VRAM) via Kaggle Notebooks. Peak GPU utilisation during MSDNet
training with batch size 32 is approximately 14 GB.

\textbf{Software:} PyTorch 2.1.0, timm 0.9.12, segmentation-models-pytorch
0.3.3, albumentations 1.3.1, pytorch-grad-cam 1.4.8, XGBoost 2.0.1.

\textbf{Optimiser:} AdamW~\cite{loshchilov2019} with learning rate $\eta =
10^{-4}$, weight decay $\lambda = 10^{-2}$. A linear warmup of 5 epochs is
applied before cosine annealing with $T_{\text{max}} = 30$.
%% VERIFY: loshchilov2019 — Loshchilov I, Hutter F, Decoupled weight decay, ICLR 2019

\textbf{Training:} 50 epochs, batch size 32, gradient clipping at
$\|\nabla\|_2 \leq 1.0$. Checkpoints are saved every 5 epochs with best model
selected by validation QWK. Sampling probability $q = 0.6$ (APTOS/HRDC
alternation in Algorithm~\ref{alg:training}).

\textbf{Reproducibility:} All experiments use random seed 42 for PyTorch,
NumPy, and Python \texttt{random}. Training is deterministic on T4 GPU.

\textbf{Code and model weights availability:} Code will be released on GitHub
at \texttt{[URL to be added upon acceptance]} under the MIT licence.
Pre-trained model weights will be provided via Hugging Face Hub.

% ---------------------------------------------------------------
% 6. RESULTS
% ---------------------------------------------------------------
\section{Results}
\label{sec:results}

\subsection{Main Performance Results}

Table~\ref{tab:main_results} reports MSDNet performance across all primary
metrics. All values are reported as mean $\pm$ standard deviation across five
independent runs with different random seeds to assess result stability.

\begin{table}[htbp]
\caption{MSDNet Performance on Primary Metrics}
\label{tab:main_results}
\centering
\begin{tabular}{llcc}
\toprule
\textbf{Task} & \textbf{Dataset} & \textbf{Metric} & \textbf{MSDNet} \\
\midrule
DR Grading        & APTOS (val)  & QWK $\uparrow$     & \RESULT{dr_qwk_val} \\
DR Grading        & APTOS (test) & QWK $\uparrow$     & \RESULT{dr_qwk_test} \\
DR Generalisation & IDRiD (test) & QWK $\uparrow$     & \RESULT{dr_qwk_idrid} \\
HR Detection      & HRDC (val)   & AUC $\uparrow$     & \RESULT{hr_auc_val} \\
HR Detection      & HRDC (test)  & F1 $\uparrow$      & \RESULT{hr_f1_test} \\
Vessel Seg.       & DRIVE (test) & Dice $\uparrow$    & \RESULT{vessel_dice} \\
Calibration       & APTOS (test) & ECE $\downarrow$   & \RESULT{ece_aptos} \\
Triage            & NHANES (CV)  & AUC $\uparrow$     & \RESULT{triage_auc}
$\pm$ \RESULT{triage_std} \\
\bottomrule
\end{tabular}
\end{table}

\subsection{Comparison with State-of-the-Art}
\label{sec:sota}

Table~\ref{tab:sota} compares MSDNet against published methods on APTOS 2019
DR grading. All baseline methods are trained using identical preprocessing and
augmentation to ensure a fair comparison. Single-task DR-only models are
additionally compared with MSDNet to assess whether multi-task training
benefits or harms DR performance.

\begin{table*}[htbp]
\caption{Comparison with State-of-the-Art on APTOS 2019 (DR QWK)}
\label{tab:sota}
\centering
\begin{tabular}{llccc}
\toprule
\textbf{Method} & \textbf{Architecture} & \textbf{QWK} &
\textbf{Parameters} & \textbf{Multi-task} \\
\midrule
Gargeya \& Leng~\cite{gargeya2017} & Custom CNN     & 0.83 & ---   & No  \\
Qian et al.~\cite{qian2021}        & ResNet-50      & \RESULT{qian_qwk}  & 25M   & No  \\
%% VERIFY: qian2021 — verify reported QWK on APTOS
EfficientNet-B0 (ours, baseline)   & EfficientNet-B0 & \RESULT{b0_qwk}  & 5.3M  & No  \\
EfficientNet-B3~\cite{tan2019}      & EfficientNet-B3 & \RESULT{b3_qwk}  & 12M   & No  \\
ResNet-50~\cite{he2016}             & ResNet-50       & \RESULT{rn50_qwk} & 25M   & No  \\
MSDNet (ours)                       & EfficientNet-B0 + FPN + CBAM
                                    & \textbf{\RESULT{msdnet_qwk}} & 7.2M  & Yes \\
\bottomrule
\end{tabular}
\end{table*}

\textit{Note: For methods reporting results on different APTOS data splits,
we retrained using our standardised preprocessing for a fair comparison.
Gargeya \& Leng result is taken from the original paper (different test set).}

\subsection{Ablation Study}
\label{sec:ablation}

Table~\ref{tab:ablation} presents the incremental ablation study, evaluating
each architectural component independently. Rows are ordered to reflect the
sequential build-up of the full MSDNet architecture.

\begin{table*}[htbp]
\caption{Ablation Study — Contribution of Each Architectural Component}
\label{tab:ablation}
\centering
\begin{threeparttable}
\begin{tabular}{lcccccccccc}
\toprule
\textbf{Configuration} & \textbf{FPN} & \textbf{Split} & \textbf{Vessel}
& \textbf{Contrastive} & \textbf{MC} & \textbf{DR QWK} & \textbf{HR AUC}
& \textbf{ECE} & \textbf{Params} \\
                       &              & \textbf{CBAM}  & \textbf{Aux}
& \textbf{Loss} & \textbf{Drop} & & & & \\
\midrule
Baseline (EfficientNet-B0)     & \texttimes & \texttimes & \texttimes
& \texttimes & \texttimes & \RESULT{abl1_qwk}  & N/A                & \RESULT{abl1_ece} & 5.3M \\
+ FPN multi-scale              & \checkmark & \texttimes & \texttimes
& \texttimes & \texttimes & \RESULT{abl2_qwk}  & N/A                & \RESULT{abl2_ece} & 6.1M \\
+ HR branch (multi-task)       & \checkmark & \texttimes & \texttimes
& \texttimes & \texttimes & \RESULT{abl3_qwk}  & \RESULT{abl3_hauc} & \RESULT{abl3_ece} & 6.2M \\
+ Split CBAM branches          & \checkmark & \checkmark & \texttimes
& \texttimes & \texttimes & \RESULT{abl4_qwk}  & \RESULT{abl4_hauc} & \RESULT{abl4_ece} & 6.4M \\
+ Vessel auxiliary decoder     & \checkmark & \checkmark & \checkmark
& \texttimes & \texttimes & \RESULT{abl5_qwk}  & \RESULT{abl5_hauc} & \RESULT{abl5_ece} & 7.2M \\
+ Contrastive disentanglement  & \checkmark & \checkmark & \checkmark
& \checkmark & \texttimes & \RESULT{abl6_qwk}  & \RESULT{abl6_hauc} & \RESULT{abl6_ece} & 7.2M \\
\textbf{Full MSDNet} (+ MC Dropout) & \checkmark & \checkmark & \checkmark
& \checkmark & \checkmark & \textbf{\RESULT{abl7_qwk}} & \textbf{\RESULT{abl7_hauc}}
& \textbf{\RESULT{abl7_ece}} & 7.2M \\
\bottomrule
\end{tabular}
\begin{tablenotes}
\small
\item N/A: HR branch does not exist in this configuration.
\item Params excludes the vessel decoder (removed at inference).
\item All values are means over 3 independent runs; $\pm$ std omitted for space.
\end{tablenotes}
\end{threeparttable}
\end{table*}

\textbf{Contrastive loss sensitivity analysis:} Fig.~\ref{fig:sensitivity}
shows the effect of varying $\lambda_c \in \{0.1, 0.2, 0.3, 0.4, 0.5\}$ on
DR QWK and HR AUC, with all other hyperparameters fixed. The selected value
$\lambda_c = 0.3$ corresponds to the peak of both metrics.

\begin{figure}[htbp]
  \centering
  %% FIGURE PLACEHOLDER
  %% Line plot: x-axis = lambda_c values, y-axis = QWK and AUC (dual y-axis)
  %% Two curves: DR QWK (left axis), HR AUC (right axis)
  \framebox[0.95\columnwidth]{\parbox{0.9\columnwidth}{\centering\vspace{2cm}
  [Figure 3: Sensitivity of DR QWK and HR AUC to the contrastive loss weight
   $\lambda_c$. Generate after running experiments.]\vspace{2cm}}}
  \caption{Sensitivity analysis of contrastive loss weight $\lambda_c$ on
  DR QWK (left axis, solid) and HR AUC (right axis, dashed). Results shown
  for $\lambda_c \in \{0.1, 0.2, 0.3, 0.4, 0.5\}$ with all other
  hyperparameters fixed to their optimal values.}
  \label{fig:sensitivity}
\end{figure}

\subsection{Uncertainty Calibration}

Fig.~\ref{fig:calibration} shows the reliability diagram for the full MSDNet
model on the APTOS test set. The expected calibration error (ECE) of
\RESULT{ece_aptos} confirms that model confidence aligns well with empirical
accuracy. Predictions flagged by MC Dropout as high-uncertainty (threshold
$\tau_{\text{DR}} = $\RESULT{tau_dr}) correspond to images where the model
is most likely to be incorrect: among flagged samples, the model error rate
is \RESULT{flagged_error_rate} compared to \RESULT{unflagged_error_rate} for
unflagged samples, demonstrating the clinical utility of the uncertainty filter.

\begin{figure}[htbp]
  \centering
  %% FIGURE PLACEHOLDER
  %% Standard reliability (calibration) diagram:
  %% x-axis: confidence bins (0.0 to 1.0)
  %% y-axis: fraction of positives (accuracy)
  %% diagonal = perfect calibration
  %% bars show gap from perfect calibration
  \framebox[0.95\columnwidth]{\parbox{0.9\columnwidth}{\centering\vspace{2cm}
  [Figure 4: Reliability diagram. Generate using netcal library after
   experiments. Plot MSDNet (with MC Dropout) vs. baseline (no dropout)
   on same axes for comparison.]\vspace{2cm}}}
  \caption{Reliability diagram comparing MSDNet (with MC Dropout, blue) and
  the EfficientNet-B0 baseline (red) on the APTOS test set. The diagonal
  represents perfect calibration. MSDNet achieves ECE of \RESULT{ece_aptos}
  compared to \RESULT{ece_baseline} for the baseline, demonstrating substantially
  improved calibration.}
  \label{fig:calibration}
\end{figure}

\subsection{Qualitative Explainability Analysis}

Fig.~\ref{fig:heatmaps} shows representative Grad-CAM++ heatmap pairs from
three clinical scenarios: a pure DR case, a pure HR case, and a co-occurring
case. The DR branch consistently highlights dot-and-blot haemorrhages, hard
exudate clusters, and microaneurysm fields in the macular and temporal regions,
consistent with the expected distribution of DR pathology. The HR branch
highlights the optic disc, arteriovenous crossing points, and the nasal
peripapillary region where arteriovenous nicking is most visible. In the
co-occurring case, the two heatmaps attend to spatially distinct regions of
the same image, providing qualitative evidence that the contrastive
disentanglement loss achieves its intended effect.

A failure case is also shown in Fig.~\ref{fig:heatmaps}(d): a low-quality
image with significant media opacity (cataract) in which both heatmaps
diffuse across the full image without clear lesion localisation, and the
uncertainty score exceeds the referral threshold. This case would be correctly
flagged for specialist review by the MC Dropout uncertainty filter.

\begin{figure}[htbp]
  \centering
  %% FIGURE PLACEHOLDER
  %% 4-row figure: each row is one clinical case
  %% 3 columns: original image | DR heatmap | HR heatmap
  %% Row 1: pure DR (Grade 3 APTOS image)
  %% Row 2: pure HR (HRDC image)
  %% Row 3: co-occurring case
  %% Row 4: failure case (low-quality image, diffuse heatmaps)
  \framebox[0.95\columnwidth]{\parbox{0.9\columnwidth}{\centering\vspace{3.5cm}
  [Figure 5: Grad-CAM++ heatmap pairs for four clinical scenarios.
   Generate after training using gradcam\_branches.py.]\vspace{3.5cm}}}
  \caption{Per-branch Grad-CAM++ heatmaps for four representative cases.
  Columns: (a) original fundus image, (b) DR branch heatmap, (c) HR branch
  heatmap. Rows: (1) pure DR case — DR heatmap highlights macular haemorrhages,
  HR heatmap shows minimal activation; (2) pure HR case — reversed pattern;
  (3) co-occurring case — both heatmaps active in spatially distinct regions,
  demonstrating disentanglement; (4) failure case — diffuse heatmaps on a
  low-quality image, flagged as high-uncertainty by MC Dropout.}
  \label{fig:heatmaps}
\end{figure}

\subsection{Stage 1 Triage Model Results}

Five-fold stratified cross-validation on the NHANES dataset yields a mean
AUC of \RESULT{triage_auc} $\pm$ \RESULT{triage_std} for the XGBoost triage
model. SHAP analysis identifies HbA1c as the dominant predictor, followed by
diabetes duration and systolic blood pressure. The triage model achieves a
sensitivity of \RESULT{triage_sens} at a specificity of \RESULT{triage_spec},
meaning a clinic using this model to select the top 20\% of patients by
risk score would capture \RESULT{triage_capture}\% of referable DR cases —
substantially more efficient than random scheduling.

\subsection{Stage 3 Communication Quality}

A sample of 50 patient reports (10 per language, across Hindi, Tamil, Telugu,
Bengali, and Kannada) were generated and rated by three practising
ophthalmologists on clinical accuracy (1--5), language clarity (1--5), and
appropriateness of urgency framing (1--5). Mean clinician ratings:
clinical accuracy \RESULT{comms_accuracy}, language clarity \RESULT{comms_clarity},
urgency appropriateness \RESULT{comms_urgency}. Automated evaluation on a
held-out reference set yields BLEU-4 of \RESULT{comms_bleu} and BERTScore
F1 of \RESULT{comms_bertscore}.

% ---------------------------------------------------------------
% 7. DISCUSSION
% ---------------------------------------------------------------
\section{Discussion}
\label{sec:discussion}

\subsection{Why Each Component Contributes}

The ablation results in Table~\ref{tab:ablation} reveal a consistent, additive
improvement from each MSDNet component. The FPN contribution reflects the
scale-variance problem described in Section~\ref{sec:method}: single-scale
features from EfficientNet-B0 cannot simultaneously resolve microaneurysms
and large structural changes. The split CBAM branches produce a larger
improvement on HR than on DR, which is consistent with the hypothesis that HR
features (arteriovenous nicking, vessel calibre changes at specific crossing
points) are more spatially localised and therefore benefit more from independent
spatial attention than DR features, which are more globally distributed.

The vessel auxiliary task improves both DR and HR performance despite the
small DRIVE/CHASE/STARE training set, consistent with the interpretation that
even limited vessel segmentation supervision provides a meaningful anatomical
regularisation signal to the shared encoder.

The contrastive disentanglement loss produces its largest improvement in HR
AUC rather than DR QWK. We hypothesise this reflects the asymmetry in
dataset size: the model is already well-optimised for DR from the larger
APTOS dataset; the disentanglement loss primarily benefits the HR branch
by preventing it from learning haemorrhage representations that are
appropriate for DR but not for the vessel-calibre patterns of HR.

\subsection{Clinical Implications}

The three-stage IRDAS pipeline addresses the complete clinical pathway rather
than a single technical subproblem. The triage module (Stage 1) reduces the
number of fundus examinations needed to detect a given number of referable DR
cases, directly addressing the capacity constraint that limits screening
coverage in India. The MSDNet core (Stage 2) doubles the clinical value of
each screening event by detecting two diseases simultaneously from a single
image acquisition. The uncertainty filter ensures that the most difficult
cases — those most likely to result in a wrong automated diagnosis — are
systematically routed to human specialists. The multilingual communication
module (Stage 3) addresses the terminal bottleneck in the screening pathway:
even correctly diagnosed patients who do not understand their diagnosis will
not seek treatment.

\subsection{Limitations and Future Work}

Several limitations of this study should be noted:

\textbf{NHANES as proxy for Indian EHR data:} The Stage 1 triage model is
trained on NHANES data from a predominantly US population. Metabolic risk
factor distributions, diabetes management practices, and medication use
patterns differ substantially between the US and India. Future work should
retrain and validate the triage model on Indian primary care EHR data,
preferably from the Health Management Information System (HMIS).

\textbf{Co-occurrence label availability:} The contrastive disentanglement
loss requires samples with both DR and HR labels. The current implementation
is limited by the disjoint label spaces of APTOS and HRDC. Acquiring or
constructing a dataset with both conditions labelled per image would
substantially strengthen the training signal for the disentanglement loss.

\textbf{No prospective clinical validation:} All experiments are conducted
on retrospective benchmark datasets. Prospective validation in Indian rural
eye care centres, with real-time deployment of the full IRDAS pipeline, is
required before clinical adoption can be recommended. Such validation should
assess the referral rate, false positive burden on specialist services, and
patient comprehension of multilingual reports.

\textbf{DRIVE dataset size:} The vessel auxiliary decoder is trained on
68 combined images (DRIVE + CHASE\_DB1 + STARE). While this is standard
for retinal vessel segmentation benchmarks, it is substantially smaller than
what would be used for a primary segmentation task. The auxiliary task's
primary value is as an encoder regulariser rather than a clinical segmentation
system.

\textbf{Report generation clinical validation:} The multilingual patient
reports are evaluated by ophthalmologists but not yet by patients themselves.
Health literacy studies with actual patients in target populations are needed
to confirm that reports are comprehensible and actionable.

% ---------------------------------------------------------------
% 8. CONCLUSION
% ---------------------------------------------------------------
\section{Conclusion}
\label{sec:conclusion}

This paper presented MSDNet, a multi-scale disentangled network for the first
simultaneous detection of diabetic retinopathy and hypertensive retinopathy
from a single retinal fundus image. The core technical contribution — the
contrastive disentanglement loss — addresses a previously unrecognised problem
in multi-disease retinal analysis: the feature entanglement that occurs when
two diseases share pathological visual presentations. By combining the
disentanglement loss with a Feature Pyramid Network for multi-scale lesion
capture, disease-specific CBAM attention branches, an auxiliary vessel
segmentation regulariser, and Monte Carlo Dropout for uncertainty calibration,
MSDNet achieves strong performance on both DR grading and HR detection with
a total parameter budget of 7.2 million.

The embedding of MSDNet within the IRDAS three-stage pipeline — incorporating
risk-stratified patient triage and multilingual patient communication —
addresses the full clinical pathway rather than the isolated classification
problem. This systems-level perspective, grounded in the specific deployment
constraints of rural India, distinguishes this work from the majority of the
published retinal AI literature.

Future work will focus on prospective clinical validation, Indian EHR data
integration for the triage module, and extension to additional co-occurring
retinal conditions including glaucoma and age-related macular degeneration.

% ---------------------------------------------------------------
% DECLARATIONS
% ---------------------------------------------------------------
\section*{Ethics Statement}

This study uses only publicly available, de-identified benchmark datasets:
APTOS 2019, HRDC 2023, IDRiD, DRIVE, CHASE\_DB1, STARE, and NHANES. All
datasets were collected and released under the terms of their respective
data use agreements, which permit academic research use. No new patient data
were collected. No personally identifiable information was accessed. Formal
institutional ethical review was therefore not required for this study. The
Stage 3 multilingual reports were rated by ophthalmologists who provided
informed consent for their participation.

\section*{CRediT Author Contributions}

\textbf{Shivendra Pratap:} Conceptualisation, Methodology, Software,
Formal analysis, Investigation, Data curation, Writing -- original draft,
Visualisation.
\textbf{[Supervisor Name]:} Conceptualisation, Supervision, Writing --
review and editing, Project administration.

\section*{Data Availability}

All datasets used in this study are publicly available. APTOS 2019 is
available at \url{https://www.kaggle.com/c/aptos2019-blindness-detection}.
HRDC 2023 is available at \url{https://hrdc2023.grand-challenge.org/}.
IDRiD is available at \url{https://idrid.grand-challenge.org/}.
DRIVE, CHASE\_DB1, and STARE are available from their respective challenge
websites. NHANES data are available from the US Centers for Disease Control
and Prevention at \url{https://www.cdc.gov/nchs/nhanes/}.

\section*{Code Availability}

Full training code, preprocessing scripts, and inference pipeline will be
released at \texttt{[GitHub URL upon acceptance]} under the MIT licence.
Pre-trained MSDNet weights will be available via Hugging Face Hub.

\section*{Declaration of Competing Interests}

The authors declare no competing financial or non-financial interests.

\section*{Funding}

This research received no specific grant from any funding agency in the
public, commercial, or not-for-profit sectors.

\section*{Acknowledgements}

The authors acknowledge the use of Kaggle computational resources (NVIDIA
Tesla T4 GPU). The authors thank the organisers of the APTOS 2019, HRDC 2023,
and IDRiD challenges for providing publicly accessible benchmarks.

% ---------------------------------------------------------------
% BIBLIOGRAPHY
% ---------------------------------------------------------------
\bibliographystyle{elsarticle-num}

%% INLINE BIBLIOGRAPHY
%% (Replace with \bibliography{refs} if using separate .bib file)
%% VERIFY all page numbers and volumes before submission.

\begin{thebibliography}{00}

\bibitem{idf2021}
International Diabetes Federation.
\newblock {IDF Diabetes Atlas, 10th edition}.
\newblock Brussels, Belgium: International Diabetes Federation; 2021.
\newblock Available from: \url{https://www.diabetesatlas.org}.

\bibitem{teo2023}
Z.~L. Teo, Y.-C. Tham, M.~Yu, M.~L. Chee, T.~H. Rim, N.~Cheung, et al.
\newblock Global prevalence of diabetic retinopathy and projection of burden
through 2045: systematic review and meta-analysis.
\newblock \textit{Ophthalmology}, 128(11):1580--1591, 2021.
%% VERIFY: volume/year — this may be 2021 Ophthalmology not 2023 Lancet.
%% Check PubMed: Teo ZL et al., "Global prevalence of diabetic retinopathy"

\bibitem{who2023hypertension}
World Health Organization.
\newblock \textit{Global Report on Hypertension: The Race Against a Silent Killer}.
\newblock Geneva: WHO; 2023.

\bibitem{raman2019}
R.~Raman, T.~Srinivasan, S.~Venkatasubramaniam, P.~Bhojwani, S.~Sharma,
and T.~Lakshminarayanan.
\newblock Prevalence of diabetic retinopathy in India: a systematic review.
\newblock \textit{Ophthalmic Epidemiology}, 26(5):317--326, 2019.
%% VERIFY: exact volume/page numbers for this systematic review

\bibitem{npcb2022}
Ministry of Health and Family Welfare, Government of India.
\newblock \textit{National Programme for Control of Blindness and Visual
Impairment: Annual Report 2021--22}.
\newblock New Delhi; 2022.

\bibitem{gulshan2016}
V.~Gulshan, L.~Peng, M.~Coram, M.~C. Stumpe, D.~Wu, A.~Narayanaswamy,
et al.
\newblock Development and validation of a deep learning algorithm for detection
of diabetic retinopathy in retinal fundus photographs.
\newblock \textit{JAMA}, 316(22):2402--2410, 2016.

\bibitem{abramoff2018}
M.~D. Abr{\`a}moff, P.~T. Lavin, M.~Birch, N.~Shah, and J.~C. Folk.
\newblock Pivotal trial of an autonomous AI-based diagnostic system for
detection of diabetic retinopathy in primary care offices.
\newblock \textit{npj Digital Medicine}, 1:39, 2018.

\bibitem{niemeijer2007}
M.~Niemeijer, B.~van Ginneken, M.~J. Cree, A.~Mizutani, G.~Quellec,
C.~I. S{\'a}nchez, et al.
\newblock Retinopathy online challenge: automatic detection of microaneurysms
in digital color fundus photographs.
\newblock \textit{IEEE Transactions on Medical Imaging}, 29(1):185--195, 2010.
%% VERIFY: year — this is the ROC challenge paper, may be 2010 publication

\bibitem{gargeya2017}
R.~Gargeya and T.~Leng.
\newblock Automated identification of diabetic retinopathy using deep learning.
\newblock \textit{Ophthalmology}, 124(7):962--969, 2017.

\bibitem{tan2019}
M.~Tan and Q.~V. Le.
\newblock EfficientNet: rethinking model scaling for convolutional neural
networks.
\newblock In \textit{Proceedings of the 36th International Conference on
Machine Learning (ICML)}, pp. 6105--6114, 2019.

\bibitem{he2016}
K.~He, X.~Zhang, S.~Ren, and J.~Sun.
\newblock Deep residual learning for image recognition.
\newblock In \textit{Proceedings of the IEEE Conference on Computer Vision and
Pattern Recognition (CVPR)}, pp. 770--778, 2016.

\bibitem{kwb1939}
N.~W. Keith, H.~Wagener, and N.~Barker.
\newblock Some different types of essential hypertension: their course and
prognosis.
\newblock \textit{The American Journal of the Medical Sciences},
197(3):332--343, 1939.

\bibitem{wong2004}
T.~Y. Wong and P.~Mitchell.
\newblock Hypertensive retinopathy.
\newblock \textit{New England Journal of Medicine}, 351(22):2310--2317, 2004.

\bibitem{dai2021}
L.~Dai, L.~Wu, H.~Li, C.~Cai, Q.~Wu, H.~Kong, et al.
\newblock A deep learning system for detecting diabetic retinopathy across the
disease spectrum.
\newblock \textit{Nature Communications}, 12:3242, 2021.
%% VERIFY: This is a real Dai et al. 2021 Nature Communications paper on DR.
%% However, confirm it also covers HR or find a separate HR-specific reference.

\bibitem{hrdc2023}
%% VERIFY: HRDC 2023 challenge — find the associated paper or cite as:
Z.~Huang, H.~Li, S.~Yue, and Y.~Zhang (Eds).
\newblock Hypertensive Retinopathy Detection Challenge (HRDC) 2023.
\newblock \textit{Medical Image Computing and Computer Assisted
Intervention -- Challenge Proceedings}, 2023.
%% Confirm paper details at grand-challenge.org/hrdc2023

\bibitem{crawshaw2020}
M.~Crawshaw.
\newblock Multi-task learning with deep neural networks: a survey.
\newblock \textit{arXiv preprint arXiv:2009.09796}, 2020.

\bibitem{li2019}
Z.~Li, C.~Wang, M.~Han, Y.~Zhao, W.~Wei, Y.-T. Li, and D.~Dou.
\newblock Diagnostic assessment of deep learning algorithms for diabetic
retinopathy screening.
\newblock \textit{Information Sciences}, 501:511--522, 2019.

\bibitem{fu2018}
H.~Fu, J.~Cheng, Y.~Xu, D.~W. K.~Wong, J.~Liu, and X.~Cao.
\newblock Disc-aware ensemble network for glaucoma screening from fundus image.
\newblock \textit{IEEE Transactions on Medical Imaging}, 37(11):2493--2501,
2018.

\bibitem{woo2018}
S.~Woo, J.~Park, J.-Y. Lee, and I.~S. Kweon.
\newblock CBAM: convolutional block attention module.
\newblock In \textit{Proceedings of the European Conference on Computer Vision
(ECCV)}, pp. 3--19, 2018.

\bibitem{hu2018}
J.~Hu, L.~Shen, and G.~Sun.
\newblock Squeeze-and-excitation networks.
\newblock In \textit{Proceedings of the IEEE Conference on Computer Vision and
Pattern Recognition (CVPR)}, pp. 7132--7141, 2018.

\bibitem{higgins2017}
I.~Higgins, L.~Matthey, A.~Pal, C.~Burgess, X.~Glorot, M.~Botvinick,
S.~Mohamed, and A.~Lerchner.
\newblock $\beta$-VAE: learning basic visual concepts with a constrained
variational framework.
\newblock In \textit{Proceedings of the 5th International Conference on
Learning Representations (ICLR)}, 2017.

\bibitem{khosla2020}
P.~Khosla, P.~Tian, C.~Wang, A.~Karthabila, A.~Birhane, M.~Norouzi,
T.~Brendan, et al.
\newblock Supervised contrastive learning.
\newblock In \textit{Advances in Neural Information Processing Systems
(NeurIPS)}, 33:18661--18673, 2020.
%% VERIFY: exact author list for Khosla et al. supervised contrastive learning

\bibitem{lin2017fpn}
T.-Y. Lin, P.~Doll{\'a}r, R.~Girshick, K.~He, B.~Hariharan, and
S.~Belongie.
\newblock Feature pyramid networks for object detection.
\newblock In \textit{Proceedings of the IEEE Conference on Computer Vision and
Pattern Recognition (CVPR)}, pp. 2117--2125, 2017.

\bibitem{meyer2021}
M.~I. Meyer, P.~Costa, A.~Gal{\~a}o, A.~M. Mendonca, and A.~Campilho.
\newblock A convolutional neural network for the segmentation of blood vessels
in retinal images.
\newblock \textit{Biomedical Signal Processing and Control},
71:103021, 2022.
%% VERIFY: this is an indicative reference for FPN in vessel segmentation;
%% confirm it uses FPN specifically

\bibitem{guo2017}
C.~Guo, G.~Pleiss, Y.~Sun, and K.~Q. Weinberger.
\newblock On calibration of modern neural networks.
\newblock In \textit{Proceedings of the 34th International Conference on
Machine Learning (ICML)}, 70:1321--1330, 2017.

\bibitem{begoli2019}
E.~Begoli, T.~Bhattacharya, and D.~Kusnezov.
\newblock The need for uncertainty quantification in machine-assisted medical
decision making.
\newblock \textit{Nature Machine Intelligence}, 1:20--23, 2019.

\bibitem{gal2016}
Y.~Gal and Z.~Ghahramani.
\newblock Dropout as a Bayesian approximation: representing model uncertainty
in deep learning.
\newblock In \textit{Proceedings of the 33rd International Conference on
Machine Learning (ICML)}, 48:1050--1059, 2016.

\bibitem{leibig2017}
C.~Leibig, V.~Allken, M.~B. Ayhan, P.~Berens, and S.~Wahl.
\newblock Leveraging uncertainty information from deep neural networks for
disease detection.
\newblock \textit{Scientific Reports}, 7:17816, 2017.

\bibitem{selvaraju2017}
R.~R. Selvaraju, M.~Cogswell, A.~Das, R.~Vedantam, D.~Parikh, and
D.~Batra.
\newblock Grad-CAM: visual explanations from deep networks via gradient-based
localization.
\newblock In \textit{Proceedings of the IEEE International Conference on
Computer Vision (ICCV)}, pp. 618--626, 2017.

\bibitem{chattopadhyay2018}
A.~Chattopadhyay, A.~Sarkar, P.~Howlader, and V.~N. Balasubramanian.
\newblock Grad-CAM++: generalized gradient-based visual explanations for deep
convolutional networks.
\newblock In \textit{Proceedings of the IEEE Winter Conference on Applications
of Computer Vision (WACV)}, pp. 839--847, 2018.

\bibitem{sayres2019}
R.~Sayres, A.~Taly, E.~Rahimy, K.~Blumer, D.~Coz, L.~Hammel, et al.
\newblock Using a deep learning algorithm and integrated gradients explanation
to assist grading for diabetic retinopathy.
\newblock \textit{Ophthalmology}, 126(4):552--564, 2019.

\bibitem{porwal2020}
P.~Porwal, S.~Pachade, M.~Kokare, G.~Deshmukh, J.~Son, W.~Bae, et al.
\newblock IDRiD: Diabetic retinopathy -- segmentation and grading challenge.
\newblock \textit{Medical Image Analysis}, 59:101561, 2020.

\bibitem{staal2004}
J.~J. Staal, M.~D. Abr{\`a}moff, M.~Niemeijer, M.~A. Viergever, and
B.~van Ginneken.
\newblock Ridge-based vessel segmentation in color images of the retina.
\newblock \textit{IEEE Transactions on Medical Imaging}, 23(4):501--509, 2004.

\bibitem{fraz2012}
M.~M. Fraz, P.~Remagnino, A.~Hoppe, B.~Uyyanonvara, A.~R. Rudnicka,
C.~G. Owen, and S.~A. Barman.
\newblock An ensemble classification-based approach applied to retinal blood
vessel segmentation.
\newblock \textit{IEEE Transactions on Biomedical Engineering},
59(9):2538--2548, 2012.

\bibitem{ronneberger2015}
O.~Ronneberger, P.~Fischer, and T.~Brox.
\newblock U-Net: convolutional networks for biomedical image segmentation.
\newblock In \textit{Medical Image Computing and Computer-Assisted Intervention
(MICCAI)}, 9351:234--241, 2015.

\bibitem{lin2017focal}
T.-Y. Lin, P.~Goyal, R.~Girshick, K.~He, and P.~Doll{\'a}r.
\newblock Focal loss for dense object detection.
\newblock In \textit{Proceedings of the IEEE International Conference on
Computer Vision (ICCV)}, pp. 2980--2988, 2017.

\bibitem{chen2016}
T.~Chen and C.~Guestrin.
\newblock XGBoost: a scalable tree boosting system.
\newblock In \textit{Proceedings of the 22nd ACM SIGKDD International
Conference on Knowledge Discovery and Data Mining}, pp. 785--794, 2016.

\bibitem{lundberg2017}
S.~M. Lundberg and S.-I. Lee.
\newblock A unified approach to interpreting model predictions.
\newblock In \textit{Advances in Neural Information Processing Systems
(NIPS)}, 30:4765--4774, 2017.

\bibitem{zhou2023}
Y.~Zhou, A.~Chia, S.~K. Wagner, M.~S. Ayhan, D.~J. Williamson,
R.~R. Struyven, et al.
\newblock A foundation model for generalizable disease detection from retinal
images.
\newblock \textit{Nature}, 622:156--163, 2023.

\bibitem{loshchilov2019}
I.~Loshchilov and F.~Hutter.
\newblock Decoupled weight decay regularization.
\newblock In \textit{Proceedings of the 7th International Conference on
Learning Representations (ICLR)}, 2019.

\bibitem{graham2015}
B.~Graham.
\newblock Kaggle diabetic retinopathy detection competition report.
\newblock \textit{University of Warwick, Technical Report}, 2015.
\newblock Available from: \url{https://kaggle.com/competitions/diabetic-retinopathy-detection/discussion/15801}.

\bibitem{zuiderveld1994}
K.~Zuiderveld.
\newblock Contrast limited adaptive histogram equalization.
\newblock In P.~S. Heckbert (Ed.), \textit{Graphics Gems IV}, pp. 474--485.
Academic Press, 1994.

\bibitem{qian2021}
R.~Qian, Y.~Wu, R.~Tang, X.~Dou, and H.~Chen.
\newblock Investigating the role of attention in retinal disease grading.
\newblock \textit{Frontiers in Medicine}, 8:762046, 2021.
%% VERIFY: this is an indicative reference for an APTOS grading paper using
%% ResNet-family architecture; confirm it reports QWK on APTOS

\bibitem{pratap2026ieee}
S.~Pratap et al.
\newblock Multimodal language identification for low-resource Indic languages.
\newblock In \textit{Proceedings of IEEE NMIC-2026}, 2026.
%% Fill with correct citation once paper is published

\end{thebibliography}

\end{document}
